\documentclass[11pt,a4paper]{article}
\usepackage{fontspec,amsmath,amssymb,amsthm,microtype}
\setmainfont{lmroman10-regular.otf}[BoldFont=lmroman10-bold.otf,ItalicFont=lmroman10-italic.otf,BoldItalicFont=lmroman10-bolditalic.otf]
\usepackage{xcolor}
\usepackage[most]{tcolorbox}
\usepackage{mathrsfs}
\providecommand{\lt}{<}
\providecommand{\gt}{>}
\setmonofont{DejaVuSansMono.ttf}[BoldFont=DejaVuSansMono-Bold.ttf,ItalicFont=DejaVuSansMono-Oblique.ttf,BoldItalicFont=DejaVuSansMono-BoldOblique.ttf,Scale=0.86]
\newfontfamily\leanfallback{STIXTwoMath-Regular.otf}[Scale=0.86]
\definecolor{leanrule}{gray}{0.62}
\definecolor{leanback}{gray}{0.962}
\definecolor{leantitle}{gray}{0.22}
\definecolor{leanmuted}{gray}{0.38}
\newtcolorbox{leanbox}[1]{enhanced,breakable,colback=leanback,colframe=leanrule,
  boxrule=0.4pt,arc=0.6mm,left=2.4mm,right=2.4mm,top=1.8mm,bottom=1.8mm,
  before skip=0.8em,after skip=1em,fonttitle=\small\bfseries,coltitle=leantitle,
  colbacktitle=leanback,titlerule=0.3pt,title={#1}}
\newlength{\leancharwidth}
\newenvironment{leancode}{\par\smallskip\begingroup\ttfamily\footnotesize\raggedright
  \settowidth{\leancharwidth}{M}\setlength{\parskip}{0pt}\setlength{\parindent}{0pt}
  \frenchspacing}{\par\endgroup\smallskip}
\newcommand{\leanline}[2]{\hangindent=\dimexpr#1\leancharwidth+4\leancharwidth\relax
  \hangafter=1\noindent\hspace*{#1\leancharwidth}#2\par}
\newcommand{\leansource}[1]{\par{\footnotesize\color{leanmuted}#1}\par}
\newcommand{\leanlabel}[1]{\par\smallskip\noindent{\small\bfseries #1}\par}
\usepackage[margin=25mm]{geometry}
\usepackage[unicode,hidelinks]{hyperref}
\hypersetup{pdfauthor={LangevinLean Project}}
\setlength{\emergencystretch}{3em}
\setlength{\parskip}{0.4em}
\setlength{\parindent}{0pt}
\allowdisplaybreaks[2]
\raggedbottom
\title{\texorpdfstring{Convolution and time-primitive identities for Wiener responses}{Convolution and time-primitive identities for Wiener responses}}
\author{LangevinLean Project}
\date{}
\begin{document}\maketitle
Draft mathematical exposition; these notes may contain mistakes.


The Gaussian innovations of the kinetic algorithms are Wiener integrals of deterministic kernels. Let \(\mathcal H=L^2(\mathbb R)\) with Lebesgue measure. For a continuous \(k:\mathbb R\to\mathbb R\) and \(t\ge0\), the time kernel \(K_k(t)\in\mathcal H\) is


\[K_k(t)(r)=\mathbf 1_{(0,t]}(r)\,k(t-r),\qquad\langle f,K_k(t)\rangle_{\mathcal H}=\int_{(0,t]}f(r)\,k(t-r)\,dr.\]


If \(I:\mathcal H\to L^2(\mathbb P_W)\) is the isonormal (Wiener) map of the Brownian motion, \(I(\mathbf 1_{(0,t]})=W_t\), and \(\tau_a\) is translation by \(a\ge0\), then \(N_k(a,t)=I(\tau_aK_k(t))=\int_a^{a+t}k(a+t-r)\,dW_r\) is the response of the kernel \(k\) to the Brownian increments on \((a,a+t]\). For instance \(\sqrt{2\gamma}\,N_{\psi_0}(a,h)\) and \(\sqrt{2\gamma}\,N_{\psi_1}(a,h)\) are the velocity and position innovations of a KLMC step, and \(N_{\varphi_2},N_{\varphi_3}\) the additional noises of KLMC2.




This section proves, for \(\gamma\ge0\) and every \(k\in\mathbb N\), the identities


\[N_{\varphi_{k+1}}(a,t)=\int_0^te^{-\gamma(t-s)}N_{\psi_k}(a,s)\,ds,\qquad N_{\psi_{k+1}}(a,t)=\int_0^tN_{\psi_k}(a,s)\,ds,\]


the second for every real \(\gamma\), as equalities of random variables, both scalar and vector-valued, on the same Brownian source. They are the stochastic Fubini identities used in the paper's derivation of KLMC2. They are obtained by proving the corresponding deterministic identities in \(\mathcal H\), with the ordinary Fubini theorem, and then applying the continuous linear map \(I\circ\tau_a\), which commutes with Bochner integrals.


\section{\texorpdfstring{The causal exponential convolution}{The causal exponential convolution}}


\textbf{Definition.} For \(\beta\in\mathbb R\) and \(k:\mathbb R\to\mathbb R\), the causal exponential convolution is


\[(\mathcal E_\beta k)(t)=\int_0^te^{-\beta(t-s)}k(s)\,ds\qquad(t\in\mathbb R),\]


an oriented integral when \(t<0\). With \(\beta=\gamma\) and \(k=\psi_j\) it is \(\varphi_{j+1}\); with \(\beta=0\) it is the primitive \(\psi_{j+1}\).




\textbf{Claim.} (a) \((\mathcal E_\beta k)(t)=e^{-\beta t}\int_0^te^{\beta s}k(s)\,ds\) for every \(t\). (b) If \(k\) is continuous, then \(\mathcal E_\beta k\) is continuous.




\textbf{Proof.} (a) For every \(s\), \(e^{-\beta(t-s)}=e^{-\beta t}e^{\beta s}\) by the addition law of the exponential, and the constant \(e^{-\beta t}\) can be taken out of the integral.




(b) The function \(s\mapsto e^{\beta s}k(s)\) is continuous, hence integrable on bounded intervals, and its primitive \(t\mapsto\int_0^te^{\beta s}k(s)\,ds\) is continuous. The product with the continuous function \(e^{-\beta t}\) is continuous, and by (a) this product is \(\mathcal E_\beta k\). \(\square\)

\begin{leanbox}{Lean: \texttt{causalExponential}, \texttt{causalExponential\_\allowbreak{}factor}, \texttt{continuous\_\allowbreak{}causalExponential}}
\leansource{\texttt{causalExponential} \quad KineticLangevin/\allowbreak{}KernelConvolution.lean, lines 11–12}
\begin{leancode}
\leanline{0}{def\ causalExponential\ (β\ :\ ℝ)\ (k\ :\ ℝ→ℝ)\ (t\ :\ ℝ)\ :\ ℝ\ :=}
\leanline{2}{∫\ s\ in\ 0..t,Real.exp\ (-β*(t-s))*k\ s}
\end{leancode}
\leansource{\texttt{causalExponential\_\allowbreak{}factor} \quad KineticLangevin/\allowbreak{}KernelConvolution.lean, lines 14–22}
\begin{leancode}
\leanline{0}{theorem\ causalExponential\_factor\ (β\ :\ ℝ)\ (k\ :\ ℝ→ℝ)\ (t\ :\ ℝ)\ :}
\leanline{4}{causalExponential\ β\ k\ t=Real.exp\ (-β*t)*∫\ s\ in\ 0..t,Real.exp\ (β*s)*k\ s}
\end{leancode}
\leansource{\texttt{continuous\_\allowbreak{}causalExponential} \quad KineticLangevin/\allowbreak{}KernelConvolution.lean, lines 24–30}
\begin{leancode}
\leanline{0}{theorem\ continuous\_causalExponential\ (β\ :\ ℝ)\ (k\ :\ ℝ→ℝ)\ (hk\ :\ Continuous\ k)\ :}
\leanline{4}{Continuous\ (causalExponential\ β\ k)}
\end{leancode}
\leanlabel{How the Lean proof follows the argument}
\texttt{causal\allowbreak{}Exponential \ensuremath{\beta} k t} is the interval integral \texttt{\ensuremath{\int} s in 0.\allowbreak{}.\allowbreak{}t,Real.\allowbreak{}exp (-\allowbreak{}\ensuremath{\beta}*(t-\allowbreak{}s))*k s}. \texttt{causal\allowbreak{}Exponential\_\allowbreak{}factor} pulls the constant inside with \texttt{interval\allowbreak{}Integral.\allowbreak{}integral\_\allowbreak{}const\_\allowbreak{}mul} and compares integrands with \texttt{interval\allowbreak{}Integral.\allowbreak{}integral\_\allowbreak{}congr}, using \texttt{Real.\allowbreak{}exp\_\allowbreak{}add}. \texttt{continuous\_\allowbreak{}causal\allowbreak{}Exponential} proves continuity of the factored form, using \texttt{interval\allowbreak{}Integral.\allowbreak{}continuous\_\allowbreak{}primitive} for the primitive of the continuous function \(e^{\beta s}k(s)\), and transfers it with \texttt{hc.\allowbreak{}congr} and \texttt{causal\allowbreak{}Exponential\_\allowbreak{}factor}.
\end{leanbox}

\section{\texorpdfstring{The two orders of integration over the triangle}{The two orders of integration over the triangle}}


The double integrals below range over the triangle \(\{(s,r):0<r\le s\le t\}\). The two identities express its sections in each order of integration.




\textbf{Claim.} Let \(t\in\mathbb R\) and \(f:\mathbb R\to\mathbb R\). (a) For any \(k\) and \(s\in(0,t]\),


\[\int_{(0,t]}\mathbf 1_{\{r\le s\}}\,k(s-r)\,f(r)\,dr=\int_{(0,s]}f(r)\,k(s-r)\,dr.\]


(b) For \(\beta,\gamma\in\mathbb R\), \(j\in\mathbb N\) and \(r\in(0,t]\),


\[\int_{(0,t]}\mathbf 1_{\{r\le s\}}\,e^{-\beta(t-s)}\psi_j(s-r)\,f(r)\,ds=f(r)\,(\mathcal E_\beta\psi_j)(t-r).\]


In both identities the two sides are integrals of the same function over the same set after the rewriting below, so they hold whenever either side is defined.




\textbf{Proof.} (a) The integrand equals the indicator of \((-\infty,s]\) times \(r\mapsto f(r)k(s-r)\). Integrating an indicator over \((0,t]\) is integrating over the intersection \((-\infty,s]\cap(0,t]\), which is \((0,s]\) because \(s\le t\).




(b) In the variable \(s\), the integrand is the indicator of \([r,\infty)\) times \(s\mapsto e^{-\beta(t-s)}\psi_j(s-r)f(r)\). Since \(r>0\), \([r,\infty)\cap(0,t]=[r,t]\), and the endpoint \(r\) has Lebesgue measure zero, so the integral is \(f(r)\int_r^te^{-\beta(t-s)}\psi_j(s-r)\,ds\). The substitution \(u=s-r\) gives


\[\int_r^te^{-\beta(t-s)}\psi_j(s-r)\,ds=\int_0^{t-r}e^{-\beta((t-r)-u)}\psi_j(u)\,du=(\mathcal E_\beta\psi_j)(t-r).\qquad\square\]\begin{leanbox}{Lean: \texttt{cutoff\_\allowbreak{}integrand\_\allowbreak{}left}, \texttt{damped\_\allowbreak{}integrand\_\allowbreak{}right}}
\leansource{\texttt{cutoff\_\allowbreak{}integrand\_\allowbreak{}left} \quad KineticLangevin/\allowbreak{}KernelConvolution.lean, lines 32–45}
\begin{leancode}
\leanline{0}{lemma\ cutoff\_integrand\_left\ (k\ :\ ℝ→ℝ)\ (t\ s\ :\ ℝ)\ (hs\ :\ s∈Ioc\ 0\ t)\ (f\ :\ ℝ→ℝ)\ :}
\leanline{4}{(∫\ r\ in\ Ioc\ 0\ t,(if\ r≤s\ then\ k\ (s-r)\ else\ 0)*f\ r)=∫\ r\ in\ Ioc\ 0\ s,f\ r*k\ (s-r)}
\end{leancode}
\leansource{\texttt{damped\_\allowbreak{}integrand\_\allowbreak{}right} \quad KineticLangevin/\allowbreak{}KernelConvolution.lean, lines 47–74}
\begin{leancode}
\leanline{0}{lemma\ damped\_integrand\_right\ (β\ γ\ :\ ℝ)\ (k\ :\ ℕ)\ (t\ r\ :\ ℝ)\ (hr\ :\ r∈Ioc\ 0\ t)\ (f\ :\ ℝ→ℝ)\ :}
\leanline{4}{(∫\ s\ in\ Ioc\ 0\ t,(if\ r≤s\ then\ Real.exp\ (-β*(t-s))*psi\ γ\ k\ (s-r)\ else\ 0)*f\ r)=}
\leanline{6}{f\ r*causalExponential\ β\ (psi\ γ\ k)\ (t-r)}
\end{leancode}
\leanlabel{How the Lean proof follows the argument}
\texttt{cutoff\_\allowbreak{}integrand\_\allowbreak{}left} rewrites the integrand as \texttt{(Iic s).\allowbreak{}indicator .\allowbreak{}.\allowbreak{}.\allowbreak{}} (\texttt{he}, by cases on \(r\le s\)), then uses \texttt{integral\_\allowbreak{}indicator} and \texttt{Measure.\allowbreak{}restrict\_\allowbreak{}restrict}; \texttt{hset} proves \((-\infty,s]\cap(0,t]=(0,s]\) by extensionality.




\texttt{damped\_\allowbreak{}integrand\_\allowbreak{}right} does the same with \texttt{Ici r}; \texttt{hset} proves \([r,\infty)\cap(0,t]=[r,t]\). Then \texttt{integral\_\allowbreak{}Icc\_\allowbreak{}eq\_\allowbreak{}integral\_\allowbreak{}Ioc} discards the endpoint, \texttt{interval\allowbreak{}Integral.\allowbreak{}integral\_\allowbreak{}of\_\allowbreak{}le} converts to an oriented integral, and \texttt{interval\allowbreak{}Integral.\allowbreak{}integral\_\allowbreak{}mul\_\allowbreak{}const} extracts \(f(r)\). \texttt{hi} is the substitution: \texttt{hfun} writes the integrand as a composition with \(s\mapsto s-r\), and \texttt{interval\allowbreak{}Integral.\allowbreak{}integral\_\allowbreak{}comp\_\allowbreak{}sub\_\allowbreak{}right} shifts the interval \([r,t]\) to \([0,t-r]\).
\end{leanbox}

\section{\texorpdfstring{The deterministic Fubini identity in \(L^2(\mathbb R)\)}{The deterministic Fubini identity in L²(R)}}


\textbf{Claim.} Let \(\beta\ge0\), \(\gamma\in\mathbb R\), \(j\in\mathbb N\) and \(t\ge0\). Then, in \(\mathcal H=L^2(\mathbb R)\),


\[K_{\mathcal E_\beta\psi_j}(t)=\int_0^te^{-\beta(t-s)}\,K_{\psi_j}(s)\,ds,\]


where the right side is the Bochner integral of the continuous \(\mathcal H\)-valued function \(s\mapsto e^{-\beta(t-s)}K_{\psi_j}(s)\).




\textbf{Proof.} Two elements of a Hilbert space are equal if they have the same inner product with every \(f\in\mathcal H\). Fix \(f\in\mathcal H\) and let \(\mu\) be Lebesgue measure on \((0,t]\).




\emph{An integrable function on the square.} Define on \((0,t]^2\)


\[w(s,r)=\begin{cases}e^{-\beta(t-s)}\psi_j(s-r)&\text{if }0<r\le s\le t,\\0&\text{otherwise.}\end{cases}\]


The function \(\psi_j\) is continuous, so it is bounded on the compact interval \([0,t]\), say \(|\psi_j|\le C\) with \(C\ge0\). On the triangle, \(s\le t\) and \(\beta\ge0\) give \(e^{-\beta(t-s)}\le1\), and \(s-r\in[0,t]\), so \(|w|\le C\) everywhere; \(w\) is measurable as a piecewise definition on a measurable set. Since \(f\in L^2\) and \((0,t]\) has finite measure, \(f\) is integrable on \((0,t]\) by the Cauchy–Schwarz inequality. Hence \(F(s,r)=w(s,r)f(r)\) is integrable for \(\mu\otimes\mu\): it is a bounded measurable function times an integrable function of \(r\). By Fubini's theorem,


\[\int_{(0,t]}\int_{(0,t]}F(s,r)\,dr\,ds=\int_{(0,t]}\int_{(0,t]}F(s,r)\,ds\,dr.\qquad(1)\]


\emph{Integrating first in \(r\).} For \(s\in(0,t]\), the previous step (a) gives


\[\int_{(0,t]}F(s,r)\,dr=e^{-\beta(t-s)}\int_{(0,s]}f(r)\psi_j(s-r)\,dr=e^{-\beta(t-s)}\langle f,K_{\psi_j}(s)\rangle.\]


\emph{Integrating first in \(s\).} For \(r\in(0,t]\), the previous step (b) gives \(\int_{(0,t]}F(s,r)\,ds=f(r)(\mathcal E_\beta\psi_j)(t-r)\), and integrating over \(r\in(0,t]\) gives \(\langle f,K_{\mathcal E_\beta\psi_j}(t)\rangle\).




\emph{Conclusion.} By (1), \(\langle f,K_{\mathcal E_\beta\psi_j}(t)\rangle=\int_0^te^{-\beta(t-s)}\langle f,K_{\psi_j}(s)\rangle\,ds\). The map \(g\mapsto\langle f,g\rangle\) is a continuous linear functional on \(\mathcal H\), so it commutes with the Bochner integral of the continuous function \(s\mapsto e^{-\beta(t-s)}K_{\psi_j}(s)\); the right side is therefore \(\langle f,\int_0^te^{-\beta(t-s)}K_{\psi_j}(s)\,ds\rangle\). Since \(f\) was arbitrary, the claim follows. \(\square\)




The hypothesis \(\beta\ge0\) is used only to bound the damping factor by \(1\) on the triangle, which yields the integrable majorant for Fubini's theorem.

\begin{leanbox}{Lean: \texttt{timeKernelLp\_\allowbreak{}causalExponential}}
\leansource{\texttt{timeKernelLp\_\allowbreak{}causalExponential} \quad KineticLangevin/\allowbreak{}KernelConvolution.lean, lines 79–153}
\begin{leancode}
\leanline{0}{theorem\ timeKernelLp\_causalExponential\ (β\ γ\ :\ ℝ)\ (hβ\ :\ 0≤β)\ (k\ :\ ℕ)\ (t\ :\ ℝ≥0)\ :}
\leanline{4}{timeKernelLp\ (causalExponential\ β\ (psi\ γ\ k))\ (continuous\_causalExponential\ β\ (psi\ γ\ k)\ (continuous\_psi\ γ\ k))\ t\ =}
\leanline{6}{∫\ s\ in\ 0..(t:ℝ),Real.exp\ (-β*((t:ℝ)-s))\ •\ timeKernelLp\ (psi\ γ\ k)\ (continuous\_psi\ γ\ k)\ s.toNNReal}
\end{leancode}
\leanlabel{How the Lean proof follows the argument}
In \texttt{time\allowbreak{}Kernel\allowbreak{}Lp\_\allowbreak{}causal\allowbreak{}Exponential}, \texttt{time\allowbreak{}Kernel\allowbreak{}Lp k hk t} is the class of \(K_k(t)\) in \texttt{Time\allowbreak{}Hilbert} (the space \(L^2(\mathbb R)\)), and \texttt{s.\allowbreak{}to\allowbreak{}NNReal} is \(\max(s,0)\), which equals \(s\) on the integration range. The proof starts with \texttt{ext\_\allowbreak{}inner\_\allowbreak{}left \ensuremath{\mathbb R}}, which reduces the equality to equality of inner products with every \texttt{f}.



\begin{enumerate}


\item \texttt{\ensuremath{\mu}} is Lebesgue measure restricted to \texttt{Ioc 0 t} and \texttt{w} is the function \(w\) above. \texttt{C}, \texttt{h\allowbreak{}C} come from \texttt{exists\_\allowbreak{}bound\_\allowbreak{}of\_\allowbreak{}continuous\allowbreak{}On} on the compact interval \texttt{Icc 0 t}, and \texttt{h\allowbreak{}C0} is \(C\ge0\).

\item \texttt{hf} is integrability of \(f\) on \((0,t]\), from \texttt{Lp.\allowbreak{}mem\allowbreak{}Lp} restricted to a finite-measure set. \texttt{hw} is measurability of \texttt{w}, via \texttt{Measurable.\allowbreak{}ite}; \texttt{hwb} is the bound \(|w|\le C\), with \texttt{hexp} the inequality \(e^{-\beta(t-s)}\le1\) from \texttt{Real.\allowbreak{}exp\_\allowbreak{}le\_\allowbreak{}one\_\allowbreak{}iff} and \texttt{h\ensuremath{\beta}}.

\item \texttt{h\allowbreak{}F} is integrability of \(F\) on the product, by \texttt{bdd\_\allowbreak{}mul}; \texttt{hswap} is identity (1), Mathlib's \texttt{integral\_\allowbreak{}integral\_\allowbreak{}swap}.

\item \texttt{hleft} evaluates the inner \(r\)-integral with \texttt{cutoff\_\allowbreak{}integrand\_\allowbreak{}left} and \texttt{inner\_\allowbreak{}time\allowbreak{}Kernel\allowbreak{}Lp} (the formula for \(\langle f,K_k(s)\rangle\)); \texttt{hright} evaluates the inner \(s\)-integral with \texttt{damped\_\allowbreak{}integrand\_\allowbreak{}right} and then the outer integral with \texttt{inner\_\allowbreak{}time\allowbreak{}Kernel\allowbreak{}Lp}.

\item \texttt{he} chains \texttt{hright}, \texttt{hswap} and \texttt{hleft}. \texttt{hc} is continuity of the integrand (using \texttt{continuous\_\allowbreak{}time\allowbreak{}Kernel\allowbreak{}Lp}), and the last line moves the inner product inside the integral with \texttt{(inner\allowbreak{}SL \ensuremath{\mathbb R} f).\allowbreak{}interval\allowbreak{}Integral\_\allowbreak{}comp\_\allowbreak{}comm}, where \texttt{inner\allowbreak{}SL \ensuremath{\mathbb R} f} is the continuous linear functional \(g\mapsto\langle f,g\rangle\).

\end{enumerate}
\end{leanbox}

\section{\texorpdfstring{The \(\varphi\) kernels and the primitives of the \(\psi\) kernels}{The φ kernels and the primitives of the ψ kernels}}


\textbf{Claim.} Let \(j\in\mathbb N\) and \(t\ge0\). (a) If \(\gamma\ge0\), then \(K_{\varphi_{j+1}}(t)=\int_0^te^{-\gamma(t-s)}K_{\psi_j}(s)\,ds\). (b) For every \(\gamma\in\mathbb R\), \(K_{\psi_{j+1}}(t)=\int_0^tK_{\psi_j}(s)\,ds\).




\textbf{Proof.} (a) By definition \(\varphi_{j+1}(u)=\int_0^ue^{-\gamma(u-s)}\psi_j(s)\,ds=(\mathcal E_\gamma\psi_j)(u)\), so this is the previous step with \(\beta=\gamma\ge0\).




(b) With \(\beta=0\), \((\mathcal E_0\psi_j)(u)=\int_0^u\psi_j(s)\,ds=\psi_{j+1}(u)\) and \(e^{-0\cdot(t-s)}=1\). The previous step with \(\beta=0\), which is allowed for every \(\gamma\), gives the claim. \(\square\)




Integration in time raises the index of \(\psi\), and damped integration turns \(\psi_j\) into \(\varphi_{j+1}\). In particular \(K_{\varphi_2}(t)=\int_0^te^{-\gamma(t-s)}\int_0^sK_{\psi_0}(r)\,dr\,ds\).

\begin{leanbox}{Lean: \texttt{timeKernelLp\_\allowbreak{}phi}, \texttt{timeKernelLp\_\allowbreak{}primitive\_\allowbreak{}all}}
\leansource{\texttt{timeKernelLp\_\allowbreak{}phi} \quad KineticLangevin/\allowbreak{}KernelConvolution.lean, lines 156–159}
\begin{leancode}
\leanline{0}{theorem\ timeKernelLp\_phi\ (γ\ :\ ℝ)\ (hγ\ :\ 0≤γ)\ (k\ :\ ℕ)\ (t\ :\ ℝ≥0)\ :}
\leanline{4}{timeKernelLp\ (phi\ γ\ (k+1))\ (continuous\_phi\ γ\ (k+1))\ t\ =}
\leanline{6}{∫\ s\ in\ 0..(t:ℝ),Real.exp\ (-γ*((t:ℝ)-s))\ •\ timeKernelLp\ (psi\ γ\ k)\ (continuous\_psi\ γ\ k)\ s.toNNReal}
\end{leancode}
\leansource{\texttt{timeKernelLp\_\allowbreak{}primitive\_\allowbreak{}all} \quad KineticLangevin/\allowbreak{}KernelConvolution.lean, lines 162–169}
\begin{leancode}
\leanline{0}{theorem\ timeKernelLp\_primitive\_all\ (γ\ :\ ℝ)\ (k\ :\ ℕ)\ (t\ :\ ℝ≥0)\ :}
\leanline{4}{timeKernelLp\ (psi\ γ\ (k+1))\ (continuous\_psi\ γ\ (k+1))\ t=}
\leanline{6}{∫\ s\ in\ 0..(t:ℝ),timeKernelLp\ (psi\ γ\ k)\ (continuous\_psi\ γ\ k)\ s.toNNReal}
\end{leancode}
\leanlabel{How the Lean proof follows the argument}
\texttt{time\allowbreak{}Kernel\allowbreak{}Lp\_\allowbreak{}phi} is the term \texttt{time\allowbreak{}Kernel\allowbreak{}Lp\_\allowbreak{}causal\allowbreak{}Exponential \ensuremath{\gamma} \ensuremath{\gamma} h\ensuremath{\gamma} k t}: the definition of \texttt{phi \ensuremath{\gamma} (k+1)} is literally \texttt{causal\allowbreak{}Exponential \ensuremath{\gamma} (psi \ensuremath{\gamma} k)}. In \texttt{time\allowbreak{}Kernel\allowbreak{}Lp\_\allowbreak{}primitive\_\allowbreak{}all}, \texttt{he} proves \texttt{causal\allowbreak{}Exponential 0 (psi \ensuremath{\gamma} k)=psi \ensuremath{\gamma} (k+1)} by unfolding (\texttt{psi\_\allowbreak{}succ}), \texttt{hi} is the general identity at \(\beta=0\) (\texttt{le\_\allowbreak{}rfl} proves \(0\le0\)), and \texttt{simpa} removes the factor \texttt{Real.\allowbreak{}exp 0 = 1}.
\end{leanbox}

\section{\texorpdfstring{Transfer to the Brownian responses, scalar and vector}{Transfer to the Brownian responses, scalar and vector}}


\textbf{Setting.} Let \(I:\mathcal H\to L^2(\mathbb P_W)\) be the isonormal map defining the scalar Brownian motion: a linear isometry with \(I(\mathbf 1_{(0,t]})=W_t\), so that \(I(g)=\int g(r)\,dW_r\) for deterministic \(g\in\mathcal H\). Let \(\tau_a g(r)=g(r-a)\), a linear isometry of \(\mathcal H\). For a continuous kernel \(k\) and \(a,t\ge0\), the scalar response is \(N_k(a,t)=I(\tau_aK_k(t))\). For \(p\) coordinates, the vector response \(\mathbf N_k(a,t)\in L^2(\mathbb P_W^{\otimes p};\mathbb R^p)\) is obtained by the spatial lift \(S\), which evaluates a scalar \(L^2\) variable on the \(i\)-th independent copy of the Brownian source and places it in the \(i\)-th coordinate: \(S(X)(\omega)=\sum_iX(\omega_i)e_i\). Both \(I\circ\tau_a\) and \(S\) are continuous linear maps.




\textbf{Claim.} For \(j\in\mathbb N\) and \(a,t\ge0\):


\[N_{\varphi_{j+1}}(a,t)=\int_0^te^{-\gamma(t-s)}N_{\psi_j}(a,s)\,ds\quad(\gamma\ge0),\qquad N_{\psi_{j+1}}(a,t)=\int_0^tN_{\psi_j}(a,s)\,ds\quad(\gamma\in\mathbb R),\]


and the same two identities hold for the vector responses \(\mathbf N\). The integrals are Bochner integrals of continuous paths in \(L^2\).




\textbf{Proof.} Let \(L=I\circ\tau_a\). By definition \(N_k(a,t)=L(K_k(t))\). Apply \(L\) to the identities of the previous step. A continuous linear map commutes with the Bochner integral of an integrable function, here the continuous integrand \(s\mapsto e^{-\gamma(t-s)}K_{\psi_j}(s)\) (respectively \(s\mapsto K_{\psi_j}(s)\)) on \([0,t]\); and it commutes with multiplication by the scalar \(e^{-\gamma(t-s)}\). Hence


\[L\bigl(K_{\varphi_{j+1}}(t)\bigr)=\int_0^tL\bigl(e^{-\gamma(t-s)}K_{\psi_j}(s)\bigr)\,ds=\int_0^te^{-\gamma(t-s)}N_{\psi_j}(a,s)\,ds,\]


and similarly for the primitive identity. For the vector responses, \(\mathbf N_k(a,t)=S(N_k(a,t))\); the scalar paths \(s\mapsto N_{\psi_j}(a,s)\) are continuous, and applying the continuous linear map \(S\) to the scalar identities gives the vector identities by the same commutation. \(\square\)




All responses are built from the same Brownian increments on \((a,a+t]\), so the identities are equalities of random vectors on one probability space, not only of their laws. With \(j=1\) and \(j=0\) they give \(\mathbf N_{\varphi_2}(a,t)=\int_0^te^{-\gamma(t-s)}\int_0^s\mathbf N_{\psi_0}(a,r)\,dr\,ds\), the identity behind the paper's formula \(\int_0^te^{-\gamma(t-s)}\int_0^s\psi_1(s-w)\,dW_w\,ds=\int_0^t\varphi_2(t-w)\,dW_w\).

\begin{leanbox}{Lean: \texttt{segmentConvolutionLp\_\allowbreak{}phi}, \texttt{segmentConvolutionLp\_\allowbreak{}primitive\_\allowbreak{}all}, \texttt{vectorSegmentLp\_\allowbreak{}phi} and 1 more}
\leansource{\texttt{segmentConvolutionLp\_\allowbreak{}phi} \quad KineticLangevin/\allowbreak{}KernelConvolution.lean, lines 172–186}
\begin{leancode}
\leanline{0}{theorem\ segmentConvolutionLp\_phi\ (γ\ :\ ℝ)\ (hγ\ :\ 0≤γ)\ (k\ :\ ℕ)\ (a\ t\ :\ ℝ≥0)\ :}
\leanline{4}{segmentConvolutionLp\ (phi\ γ\ (k+1))\ (continuous\_phi\ γ\ (k+1))\ a\ t=}
\leanline{6}{∫\ s\ in\ 0..(t:ℝ),Real.exp\ (-γ*((t:ℝ)-s))\ •}
\leanline{8}{segmentConvolutionLp\ (psi\ γ\ k)\ (continuous\_psi\ γ\ k)\ a\ s.toNNReal}
\end{leancode}
\leansource{\texttt{segmentConvolutionLp\_\allowbreak{}primitive\_\allowbreak{}all} \quad KineticLangevin/\allowbreak{}KernelConvolution.lean, lines 189–198}
\begin{leancode}
\leanline{0}{theorem\ segmentConvolutionLp\_primitive\_all\ (γ\ :\ ℝ)\ (k\ :\ ℕ)\ (a\ t\ :\ ℝ≥0)\ :}
\leanline{4}{segmentConvolutionLp\ (psi\ γ\ (k+1))\ (continuous\_psi\ γ\ (k+1))\ a\ t=}
\leanline{6}{∫\ s\ in\ 0..(t:ℝ),segmentConvolutionLp\ (psi\ γ\ k)\ (continuous\_psi\ γ\ k)\ a\ s.toNNReal}
\end{leancode}
\leansource{\texttt{vectorSegmentLp\_\allowbreak{}phi} \quad KineticLangevin/\allowbreak{}KernelConvolution.lean, lines 203–215}
\begin{leancode}
\leanline{0}{theorem\ vectorSegmentLp\_phi\ (γ\ :\ ℝ)\ (hγ\ :\ 0≤γ)\ (k\ :\ ℕ)\ (a\ t\ :\ ℝ≥0)\ :}
\leanline{4}{vectorSegmentLp\ ι\ (phi\ γ\ (k+1))\ (continuous\_phi\ γ\ (k+1))\ a\ t=}
\leanline{6}{∫\ s\ in\ 0..(t:ℝ),Real.exp\ (-γ*((t:ℝ)-s))\ •}
\leanline{8}{vectorSegmentLp\ ι\ (psi\ γ\ k)\ (continuous\_psi\ γ\ k)\ a\ s.toNNReal}
\end{leancode}
\leansource{\texttt{vectorSegmentLp\_\allowbreak{}primitive\_\allowbreak{}all} \quad KineticLangevin/\allowbreak{}KernelConvolution.lean, lines 218–225}
\begin{leancode}
\leanline{0}{theorem\ vectorSegmentLp\_primitive\_all\ (γ\ :\ ℝ)\ (k\ :\ ℕ)\ (a\ t\ :\ ℝ≥0)\ :}
\leanline{4}{vectorSegmentLp\ ι\ (psi\ γ\ (k+1))\ (continuous\_psi\ γ\ (k+1))\ a\ t=}
\leanline{6}{∫\ s\ in\ 0..(t:ℝ),vectorSegmentLp\ ι\ (psi\ γ\ k)\ (continuous\_psi\ γ\ k)\ a\ s.toNNReal}
\end{leancode}
\leanlabel{How the Lean proof follows the argument}
\texttt{segment\allowbreak{}Convolution\allowbreak{}Lp k hk a t} is defined as the isonormal map \texttt{ULA.\allowbreak{}Gaussian\allowbreak{}Hilbert.\allowbreak{}isonormal ULA.\allowbreak{}Canonical\allowbreak{}Brownian.\allowbreak{}basis} applied to the translated kernel \texttt{time\allowbreak{}Shift a (time\allowbreak{}Kernel\allowbreak{}Lp k hk t)}; \texttt{vector\allowbreak{}Segment\allowbreak{}Lp \ensuremath{\iota} k hk a t} is \texttt{spatial\allowbreak{}Lift \ensuremath{\iota}} applied to it.



\begin{enumerate}


\item In \texttt{segment\allowbreak{}Convolution\allowbreak{}Lp\_\allowbreak{}phi}, \texttt{L} is the composition of the isonormal map with \texttt{time\allowbreak{}Shift a}, both converted to continuous linear maps. After \texttt{change}, the goal is \(L(K_{\varphi_{j+1}}(t))=\dots\); \texttt{time\allowbreak{}Kernel\allowbreak{}Lp\_\allowbreak{}phi} rewrites the kernel, \texttt{L.\allowbreak{}interval\allowbreak{}Integral\_\allowbreak{}comp\_\allowbreak{}comm} moves \(L\) inside the integral (it needs the interval integrability of the continuous integrand \texttt{hc}), and \texttt{map\_\allowbreak{}smul L} handles the scalar factor.

\item \texttt{segment\allowbreak{}Convolution\allowbreak{}Lp\_\allowbreak{}primitive\_\allowbreak{}all} is the same with \texttt{time\allowbreak{}Kernel\allowbreak{}Lp\_\allowbreak{}primitive\_\allowbreak{}all}.

\item \texttt{vector\allowbreak{}Segment\allowbreak{}Lp\_\allowbreak{}phi} and \texttt{vector\allowbreak{}Segment\allowbreak{}Lp\_\allowbreak{}primitive\_\allowbreak{}all} unfold \texttt{vector\allowbreak{}Segment\allowbreak{}Lp}, rewrite with the scalar identities, and commute \texttt{spatial\allowbreak{}Lift \ensuremath{\iota}} with the integral by \texttt{interval\allowbreak{}Integral\_\allowbreak{}comp\_\allowbreak{}comm}, using continuity of the scalar responses (\texttt{continuous\_\allowbreak{}segment\allowbreak{}Convolution\allowbreak{}Lp}).

\end{enumerate}
\end{leanbox}


Earlier results used here: \cite[Section~3.6.12]{LangevinLean2026ULASC}; \cite[Section~3.8.5]{LangevinLean2026ULASC}.

\begin{thebibliography}{ULA-LSI}
\bibitem[ULA-SC]{LangevinLean2026ULASC}
LangevinLean Project. \emph{Unadjusted Langevin under Strong Convexity}. LL-01, 2026. Draft report, content revision d34ead06987c. \url{https://odakyildiz.github.io/LangevinLean/docs/pdf/UnadjustedLangevin/FullProof.pdf}.
\end{thebibliography}

\end{document}
